% ===== Block D: PDF p.16–20 | 原文 §2.3（前半） | 唯一编号公式 2.36–2.61 =====
% [NOTE: 本块从 p.16 页面顶部开始；原文把 §2.3 标题排在 p.16 中部（(2.38) 之后），该标题由本块就地写入（Lead 裁决）。]
% [原文 p.16]

用微分形式来说，把 $\phi$ 局部地写成 $(x_1,\dots,x_r,y_1,\dots,y_n)\mapsto(y_1,\dots,y_n)$，并设 $\omega$ 是 $M$ 上的一个 $k$-形式，它局部地表示为 $\omega=\omega^{\nu_1,\dots,\nu_{k-r}}(x,y)\,\dd x_1\wedge\cdots\wedge \dd x_r\wedge \dd y_{\nu_1}\wedge\cdots\wedge \dd y_{\nu_{k-r}}+\cdots$。省略号代表含较少 $\dd x$ 因子的那些项。于是
% [Tutor Note: $\omega$ 的分量必须是上标 $\omega^{\nu_1,\dots,\nu_{k-r}}$、而 $y$ 的微分配下标 $\dd y_{\nu_1}\wedge\cdots$，两者缩并；文字层把该上标误抄成下标，按页面图可确认是上标。次数核对：$\dd x$ 占 $r$ 次、$\dd y$ 占 $k-r$ 次，合计 $k$。]
\begin{equation}
  (\phi_*\omega)_q=\left(\int_{\phi^{-1}(q)}\omega^{\nu_1,\dots,\nu_{k-r}}(x,q)\,\dd x_1\wedge\cdots\wedge \dd x_r\right)\dd_q y_{\nu_1}\wedge\cdots\wedge \dd_q y_{\nu_{k-r}}
  \tag{2.36}
\end{equation}
对一切 $q\in N$ 成立。推前与加法相容，但不保持杯积。
% [Tutor Note: “for all $q\in N$” 中的 $N$ 是 $\phi$ 的目标流形（与 $M$ 同级的记法），不是自然数集；原文排斜体 $N$，勿与 $\mathbb{N}$ 混淆。]

只要 $M$ 与 $N$ 都紧，该定义就可以借助 Poincaré 对偶加以推广。此时，推前无非是把同调中的推前用 Poincaré 对偶作前复合与后复合：
% [Tutor Note: $\overset{\PD}{\cong}$ 是“经 Poincaré 对偶等同”的标注，不是普通的同构记号。次数核对：$\phi$ 的纤维维数 $r=m-n$，沿纤维积分把 $k$ 次类降为 $k-(m-n)$ 次，恰与 (2.36) 中积分掉 $r$ 个 $\dd x$ 因子相吻合。]
\begin{equation}
  \phi_*:H^k(M)\overset{\PD}{\cong}H_{m-k}(M)\longrightarrow H_{m-k}(N)\overset{\PD}{\cong}H^{k-(m-n)}(N).
  \tag{2.37}
\end{equation}
只要 $\phi$ 具有紧的、维数为 $r=m-n$ 的纤维，它就与\emph{沿纤维积分}一致。

\emph{投影公式。}在纤维紧的情形，有一个有用的公式，称为投影公式，它把 $M$ 上的积分表示成 $N$ 上的积分。更确切地说：若 $\omega\in H^k(M)$、$\eta\in H^{m-k}(N)$，则
% [Tutor Note: (2.38) 的次数是“错位”配对的：$\phi^*\eta$ 属于 $H^{m-k}(M)$，故 $\omega\wedge\phi^*\eta$ 是 $M$ 上的 $m$ 次顶层形式；右端 $\phi_*\omega\in H^{k-r}(N)$ 与 $\eta$ 相乘得 $k-r+m-k=m-r=n$ 次。若误把 $\eta$ 取在 $H^{n-k}(N)$ 就无法配平，这正是投影公式的要点。]
\begin{equation}
  \int_M\omega\wedge\phi^*\eta=\int_N\phi_*\omega\wedge\eta .
  \tag{2.38}
\end{equation}

\subsection{相交理论与 Witten 猜想（Intersection theory and Witten's conjecture）}

回忆一下我们的主要目标：定义并计算黎曼曲面模空间上的积分。由于 $\Mc_{g,n}$ 是紧轨形，我们现在终于可以讨论最高次上同调类的积分。然而，我们还没有任何自然的类可以积分。上同调类有两个自然的来源。
\begin{itemize}
  \item 自然（复）子空间的 Poincaré 对偶；
  \item 自然复向量丛的 Chern 类。
\end{itemize}
这两种情形产生的都是偶数次的上同调类。出于这个原因，在上同调中作乘法时，我们将始终省略帽积，因为偶数次上同调类的帽积是可交换的。
% [Tutor Note: 此处的 “cap product”（帽积）应为 “cup product”（杯积）：两个上同调类之间只有杯积，且偶数次类的杯积可交换；同页前文刚说过推前 “does not respect the cup product”，可作旁证。已登记 ERRATA。]

我们已经遇到过 $\Mc_{g,n}$ 的若干子空间，即边界层。回忆：对一个固定的 $(g,n)$ 型稳定图 $\Gamma$，相应的子空间 $\Mc_\Gamma$ 的复维数为 $\dim(\Mc_{g,n})-|E(\Gamma)|$。由此可知，其 Poincaré 对偶（记作方括号 $[\cdot]$）落在
% [Tutor Note: 次数核对：$\Mc_\Gamma$ 的复余维为 $|E(\Gamma)|$，即实余维 $2|E(\Gamma)|$；上同调类按实次数记，故 $[\Gamma]\in H^{2|E(\Gamma)|}$。取 $|E(\Gamma)|=1$（边界除子）验证：$[\Gamma]\in H^2(\Mc_{g,n},\Q)$。]
\begin{equation}
  [\Gamma]\in H^{2|E(\Gamma)|}(\Mc_{g,n},\Q).
  \tag{2.39}
\end{equation}
它可以写成沿粘合映射的推前：
% [Tutor Note: 系数 $1/|\Aut(\Gamma)|$ 是栈（轨形）意义下的归一化：粘合映射 $\xi_\Gamma$ 在栈层面是 $|\Aut(\Gamma)|$ 重覆盖，故 $\xi_{\Gamma,*}\mathbf{1}$ 须除以 $|\Aut(\Gamma)|$ 才给出正确的 Poincaré 对偶。]
\begin{equation}
  [\Gamma]=\frac{1}{|\Aut(\Gamma)|}\xi_{\Gamma,*}\mathbf{1}.
  \tag{2.40}
\end{equation}

% [原文 p.17]

右端的元素 $\mathbf{1}$ 是 $H^{\bullet}(\Mc_\Gamma,\Q)$ 中的单位元。特别地，整个空间的 Poincaré 对偶——它对应亏格为 $g$ 的单个顶点、没有边、有 $n$ 个叶子的稳定图——就是上同调中的单位元：
% [Tutor Note: 原文 (2.41) 在方括号内嵌一幅小图（$n$ 个叶子 $1,\dots,n$ 接到一个亏格 $g$ 的顶点，无其他边）；本译本不重画图，改用 $\textcircled{g}$ 表示顶点的亏格。核对：该图无边界（$|E(\Gamma)|=0$），故 $[\Gamma]\in H^0$，与右端单位元 $\mathbf{1}$ 相符。]
\begin{equation}
  \left[\;\begin{array}{c}1\\ \vdots\\ n\end{array}\;\;\mbox{\textcircled{$g$}}\;\right]=\mathbf{1}\in H^0(\Mc_{g,n},\Q).
  \tag{2.41}
\end{equation}

现在我们讨论复向量丛的 Chern 类。$\Mc_{g,n}$ 上的一个复向量丛，就是对每个稳定黎曼曲面的同构类指派一个复向量空间，使得当稳定黎曼曲面在模空间中变动时，所派定的向量空间光滑地变动并协调地粘合在一起。给定复向量丛 $\cV\to\Mc_{g,n}$ 之后，我们就可以考虑它的 Chern 类：
\begin{equation}
  c_k(\cV)\in H^{2k}(\Mc_{g,n},\Q),\qquad k=0,1,\dots,\operatorname{rk}(\cV),
  \tag{2.42}
\end{equation}
其中 $\operatorname{rk}(\cV)$ 表示 $\cV$ 的复秩，即纤维的复维数。第零个 Chern 类总是上同调中的单位元：$c_0(\cV)=\mathbf{1}$。Chern 类是复向量丛的拓扑不变量，并给出一个判定两个向量丛不同构的简单检验：若一对向量丛的 Chern 类不同，则这两个向量丛不同（反之则未必成立）。几何上，它们提供关于向量丛有多少个线性无关截面的信息，并且可以表示为 $\cV$ 上任一 Hermitian 联络 $\nabla$ 的曲率形式的特征多项式的系数（该上同调类不依赖于联络的选取）：
% [Tutor Note: (2.43) 中 $\mathrm{Id}$ 是单位矩阵、$\ii$ 是虚数单位（原文皆排正体），故不能写成斜体 $i$。核对：右端行列式关于 $t$ 的常数项为 $\det(\mathrm{Id})=1$，恰等于 $c_0(\cV)$，与正文“第零个 Chern 类总是单位元”一致。]
\begin{equation}
  c(\cV;t)=\sum_{k=0}^{\operatorname{rk}(\cV)}c_k(\cV)\,t^k=\det\left(\mathrm{Id}-t\frac{F_\nabla}{2\pi\ii}\right).
  \tag{2.43}
\end{equation}

这种全纯向量丛的第一个例子是所谓的\emph{第 $i$ 个余切线丛}：对每个 $i\in\{1,\dots,n\}$，令
\begin{equation}
  \cL_i\longrightarrow\Mc_{g,n},\qquad \cL_i|_{(\Sigma,p_1,\dots,p_n)}=T^*_{p_i}\Sigma .
  \tag{2.44}
\end{equation}
% [Tutor Note: $\cL_i|_{(\Sigma,p_1,\dots,p_n)}$ 是“点 $(\Sigma,p_1,\dots,p_n)$ 处的纤维”的记号，不是函数的限制；纤维是余切空间 $T^*_{p_i}\Sigma$，故 $\cL_i$ 是线丛（复秩 $1$）。]

换言之，$(\Sigma,p_1,\dots,p_n)$ 上的纤维是第 $i$ 个标记点处的全纯余切空间。由于 $T^*_{p_i}\Sigma$ 是复维数为 $1$ 的复向量空间，相应的丛 $\cL_i$ 的复秩为 $1$：它是一个线丛。于是我们考虑它的第一个 Chern 类：
\begin{equation}
  \psi_i=c_1(\cL_i)\in H^2(\Mc_{g,n},\Q).
  \tag{2.45}
\end{equation}

这些类称为余切线类，或简称 $\psi$ 类。按惯例，严格说来 $\psi$ 类应当依赖于 $(g,n)$。我们省略这一依赖，但愿这在上下文中是清楚的。我们很快会看到，$\psi$ 类出现在 Witten 关于拓扑二维引力的开创性工作 [Wit91] 中，并且构成联系到黎曼曲面模空间的一切物理理论（例如 JT 引力与拓扑弦论）的基石。
% [Tutor Note: 原文在此声明省略 $\psi_i$ 对 $(g,n)$ 的依赖；读后面的公式时须按所在模空间理解同一个符号，例如 (2.54) 两端的 $\psi_i$ 分别属于 $\Mc_{g,n}$ 与 $\Mc_{g,n+1}$，正是原文紧接 (2.54) 的那句话要说明的。]

% [原文 p.18]

\figomit{12}{余切线丛 $\cL_i$ 的图示。}
% [Tutor Note: Figure 12 画的是 $\cL_i\to\Mc_{g,n}$：下方的曲面示意 $\Mc_{g,n}$ 中的点 $(\Sigma,p_1,\dots,p_n)$，其上的纤维是第 $i$ 个标记点处的余切线 $T^*_{p_i}\Sigma$（图中在原图里直接标注了这个记号）。]

由 $\psi$ 类我们可以导出新的上同调类，它们是被遗忘点的投影：Morita–Miller–Mumford 类，或简称 $\kappa$ 类，定义为
% [Tutor Note: 次数核对：$\psi_{n+1}^{m+1}$ 是 $2(m+1)$ 次，$\pi_*$ 沿一维纤维积分降 $2$ 次，得 $H^{2m}$。上界 $m\le 3g-3+n=\dim_{\C}\Mc_{g,n}$ 恰好保证次数不超过顶层次数，超出即为零类。]
\begin{equation}
  \kappa_m=\pi_*\left(\psi_{n+1}^{m+1}\right)\in H^{2m}(\Mc_{g,n},\Q),\qquad m=0,\dots,3g-3+n,
  \tag{2.46}
\end{equation}
其中 $\pi:\Mc_{g,n+1}\to\Mc_{g,n}$ 是遗忘映射。由于 $\pi$ 的纤维紧且是一维的，推前在上同调中良定义，并把复上同调次数降低 $1$。我们很快会看到，类 $2\pi^2\kappa_1$（称为\emph{Weil–Petersson 类}）在 JT 理论与双曲几何中起着基本的作用。

第三组自然的上同调类来自与黎曼曲面相联系的最自然的向量空间：全纯微分形式空间。更确切地说，定义 Hodge 丛
\begin{equation}
  \cH\longrightarrow\Mc_{g,n},\qquad \cH|_{(\Sigma,p_1,\dots,p_n)}=\Omega(\Sigma).
  \tag{2.47}
\end{equation}
这里 $\Omega(\Sigma)$ 表示 $\Sigma$ 上全纯形式构成的空间，它是复维数为 $g$ 的复向量空间。然而，对于带结点的黎曼曲面，全纯形式的定义（也就是模空间边界上 $\cH$ 的定义）需要小心。为了理解在结点黎曼曲面上应当如何定义全纯形式，考虑例子
% [Tutor Note: $\cH$ 的秩为 $g=\dim\Omega(\Sigma)$，这正是 (2.52) 中 $k$ 只到 $g$ 的原因；$\lambda_0=c_0(\cH)=1$。]
\begin{equation}
  E_t:\quad y^2=x(x-1)(x-t).
  \tag{2.48}
\end{equation}
对 $t\ne 0$，$E_t$ 上全纯形式的空间是一维的，由
\begin{equation}
  \omega_t=\frac{\dd x}{y}=\frac{\dd x}{\sqrt{x(x-1)(x-t)}}.
  \tag{2.49}
\end{equation}
生成。当 $t\to 0$ 时，环面退化成一个捏缩环面，全纯形式 $\omega_t$ 的极限为
% [Tutor Note: 由 (2.49) 取极限即可核对 (2.50)：$t\to0$ 时 $\sqrt{x(x-1)(x-t)}\to\sqrt{x^2(x-1)}$，按分支 $\sqrt{x^2}=x$ 得 $\dd x/(x\sqrt{x-1})$。$\omega_0$ 在 $x=0$ 处有一阶极点；$\sqrt{x-1}$ 的两支相差一个符号，故两支上留数相反。另可核验 (2.49)：$\omega_t=\dd x/y$ 在分支点处仍全纯（局部 $y\sim\sqrt{x-a}$，$\dd x/y\sim 2\dd y$）。]
\begin{equation}
  \omega_0=\frac{\dd x}{x\sqrt{x-1}}.
  \tag{2.50}
\end{equation}

% [原文 p.19]

在局部坐标中可以验证，$\omega_0$ 不再是全纯的，而是亚纯的：它在结点处有一阶极点，且在结点的两个分支上留数相反。这个一阶极点的出现是关键的。事实上，捏缩环面就是把 $\PP^1$ 的两个点等同起来；$\PP^1$ 上没有非平凡的全纯形式；但存在一维复向量空间，由在这两个特殊点处有一阶极点、且留数相反的亚纯形式组成。换言之，即使在极限 $t\to 0$ 之下，$\Omega(E_t)$ 的维数也保持不变。

于是 $\Omega(\Sigma)$ 的定义为
\begin{equation}
  \Omega(\Sigma)=\left\{\;
  \begin{array}{c}
    \mbox{$\Sigma$ 上的亚纯形式}\\
    \mbox{在结点处至多有一阶极点，留数相反}\\
    \mbox{而在其余各处全纯}
  \end{array}
  \;\right\},
  \tag{2.51}
\end{equation}
当 $\Sigma$ 在 $\Mc_{g,n}$ 中变动时，它的维数恒为 $g$（与标记点无关）。接着我们定义 Hodge 类，或简称 $\lambda$ 类，为 Hodge 丛的 Chern 类：
% [Tutor Note: (2.51) 的要点是“允许一阶极点、要求留数相反”：只有这样才能让 $\dim\Omega$ 在退化时保持不变。若改取全纯形式，则 (2.50) 的 $\omega_0$ 会被排除，维数就从 $g=1$ 掉到 $0$。]
\begin{equation}
  \lambda_k=c_k(\cH)\in H^{2k}(\Mc_{g,n},\Q),\qquad k=0,\dots,g,
  \tag{2.52}
\end{equation}
正如我们将在第 4 节简要提到的，Hodge 类在拓扑弦论中起着基本作用。

本节最后，我们对 Witten 猜想作一个简要的概述。我们从关于 $\psi$ 类相交数（也称为 \emph{Witten 关联函数}）的两个事实出发，即弦方程与膨胀子方程（string and dilaton equations）。这些方程把 $\psi$ 类在不同模空间上的积分联系起来。按照 Witten 的记号，这类积分可以方便地写成
% [Tutor Note: 非零条件 $d_1+\cdots+d_n=3g-3+n$ 正是“被积类是顶层类”：$\dim_{\C}\Mc_{g,n}=3g-3+n$，次数低于 $2(3g-3+n)$ 时积分为零，故 (2.53) 的定义只需考虑满足该等式的 $(d_1,\dots,d_n)$。]
\begin{equation}
  \langle\tau_{d_1}\cdots\tau_{d_n}\rangle_g=\int_{\Mc_{g,n}}\psi_1^{d_1}\cdots\psi_n^{d_n},\qquad d_i\ge 0.
  \tag{2.53}
\end{equation}
除非 $d_1+\cdots+d_n=3g-3+n$，该积分约定为零；而在该情形下，被积式是最高维的上同调类。

\begin{itemize}
  \item \emph{几何弦方程。}沿遗忘映射的 $\psi$ 类拉回为
  % [Tutor Note: (2.54) 方括号内是稳定图：左边亏格 $g$ 的顶点带标记 $1,\dots,\hat{\imath},\dots,n$（帽号表示略去 $i$），右边亏格 $0$ 的顶点带标记 $i$ 与 $n+1$，两顶点由一条边相连（图中 $|E|=1$，故 $D_i$ 是除子，即 $H^2$ 类）。本译本不重画图，用两个圆圈与其间的横线表示这条边。]
  \begin{equation}
    \pi^*\psi_i=\psi_i-D_i,\qquad
    D_i=\left[\;
    \begin{array}{c}1\\ \hat{\imath}\\ \vdots\\ n\end{array}
    \;\;\mbox{\textcircled{$g$}}\ \mbox{---}\ \mbox{\textcircled{$0$}}\;\;
    \begin{array}{c}i\\ \\ n+1\end{array}
    \;\right],
    \tag{2.54}
  \end{equation}
  其中照例，像 $\hat{\imath}$ 中的帽号表示略去。左端的 $\psi$ 类位于 $\Mc_{g,n}$ 中，而右端的则位于 $\Mc_{g,n+1}$ 中。
  \item \emph{几何膨胀子方程。}$\Mc_{g,n}$ 上的第 $0$ 个 $\kappa$ 类等于（负的）Euler 示性数：
  \begin{equation}
    \kappa_0=(2g-2+n)\mathbf{1}\in H^0(\Mc_{g,n},\Q).
    \tag{2.55}
  \end{equation}
  % [Tutor Note: 核对符号：$\chi(\Mc_{g,n})=2-2g-n$，故 (2.55) 右端 $2g-2+n=-\chi$，这正是正文“(minus) the Euler characteristic”的含义。特例 $g=0,n=3$：$\Mc_{0,3}$ 是一个点，$\kappa_0=1$，与 $2g-2+n=1$ 相符。]
\end{itemize}

\begin{exercise}[2.5]
利用几何弦方程与几何膨胀子方程，连同投影公式以及边界层的 Poincaré 对偶的表达式 (2.40)，证明 Witten 关联函数所满足的下列方程。

% [原文 p.20]

\begin{itemize}
  \item \emph{弦方程。}在 $\Mc_{g,n+1}$ 上不含 $\psi_{n+1}$ 的积分可化归为 $\Mc_{g,n}$ 上的积分：
  % [Tutor Note: (2.56) 与 (2.57) 是同一方程的两种写法：把 (2.56) 的 $\psi_i^{d_i}$ 降一次幂即得 $\tau_{d_i}\to\tau_{d_i-1}$。核对 $g=0$、$n=3$、$(d_1,d_2,d_3)=(1,0,0)$：左边 $\int_{\Mc_{0,4}}\psi_1=1$，右边只有 $i=1$ 项非零，等于 $\int_{\Mc_{0,3}}1=1$，两边一致（$i=2,3$ 项含 $\psi^{-1}$，为零）。]
  \begin{equation}
    \int_{\Mc_{g,n+1}}\psi_1^{d_1}\cdots\psi_n^{d_n}=\sum_{i=1}^{n}\int_{\Mc_{g,n}}\left(\prod_{j\ne i}\psi_j^{d_j}\right)\psi_i^{d_i-1}.
    \tag{2.56}
  \end{equation}
  按 Witten 的记号，弦方程相当于去掉一个 $\tau_0$：
  \begin{equation}
    \langle\tau_{d_1}\cdots\tau_{d_n}\tau_0\rangle_g=\sum_{i=1}^{n}\langle\tau_{d_1}\cdots\tau_{d_i-1}\cdots\tau_{d_n}\rangle_g .
    \tag{2.57}
  \end{equation}
\end{itemize}

\noindent 提示。考虑以下事实。
% [Tutor Note: 原文用一枚“指示手”花饰引出 Hints（本译本不排该花饰）。另注意原文把 Exercise 2.5 的提示一段插在两条方程之间（先弦方程、再提示、再膨胀子方程），本译本照原顺序保留；第一条提示其实就是用 (2.56) 中推前降低次数来论证。]
\begin{itemize}[label=--]
  \item 从同调次数看，对同调类 $\alpha\in H^{2(3g-3+n)}(\Mc_{g,n},\Q)$，积分 $\int_{\Mc_{g,n+1}}\pi^*\alpha$ 能说明什么？
  \item 设 $D_i$ 如 (2.54) 式中所定义。把它解释为 Poincaré 对偶，可以看出对所有 $i\ne j$ 都有 $D_i\cdot D_j=0$。
\end{itemize}

\begin{itemize}
  \item \emph{膨胀子方程。}在 $\Mc_{g,n+1}$ 上带单个 $\psi_{n+1}$ 幂次的积分可化归为 $\Mc_{g,n}$ 上的积分：
  \begin{equation}
    \int_{\Mc_{g,n+1}}\psi_1^{d_1}\cdots\psi_n^{d_n}\psi_{n+1}=(2g-2+n)\int_{\Mc_{g,n}}\psi_1^{d_1}\cdots\psi_n^{d_n}.
    \tag{2.58}
  \end{equation}
  按 Witten 的记号，膨胀子方程相当于去掉一个 $\tau_1$：
  % [Tutor Note: 次数核对：左端比右端多一个 $\tau_1$，$d$ 之和与插入点数同时加 $1$，故两边的非零条件 $d$ 之和 $=3g-3+n$ 完全相容，系数只能是一个只依赖 $g,n$ 的常数；再由 (2.55)（把 $\psi_{n+1}$ 换成 $\psi_{n+1}^0$）配合投影公式即可定出该常数为 $2g-2+n$。]
  \begin{equation}
    \langle\tau_{d_1}\cdots\tau_{d_n}\tau_1\rangle_g=(2g-2+n)\langle\tau_{d_1}\cdots\tau_{d_n}\rangle_g .
    \tag{2.59}
  \end{equation}
\end{itemize}
\end{exercise}

弦方程与膨胀子方程使得亏格 $0$ 与亏格 $1$ 的所有 Witten 关联函数都可以计算出来。

\begin{exercise}[2.6]
已知弦方程与积分 $\langle\tau_0^3\rangle_0=1$，证明所有亏格 $0$ 的 $\psi$ 类相交数都被确定。你能证明下面的闭公式吗：
% [Tutor Note: (2.60) 可用小值核对：$n=3$、$(d_1,d_2,d_3)=(0,0,0)$ 给出 $\binom{0}{0,0,0}=1=\langle\tau_0^3\rangle_0$；$n=4$、$(1,0,0,0)$ 给出 $\binom{1}{1,0,0,0}=1=\int_{\Mc_{0,4}}\psi_1$。注意多重系数的上标 $n-3$ 恰等于 $\sum d_i$（对亏格 0 的非零情形），否则该系数无数值意义。]
\begin{equation}
  \langle\tau_{d_1}\cdots\tau_{d_n}\rangle_0=\binom{n-3}{d_1,\dots,d_n},
  \tag{2.60}
\end{equation}
其中 $\binom{D}{d_1,\dots,d_n}=\frac{D!}{d_1!\cdots d_n!}$ 是多项系数？
\end{exercise}

\begin{exercise}[2.7]
已知弦方程、膨胀子方程以及积分 $\langle\tau_1\rangle_1=\frac{1}{24}$，证明所有亏格 $1$ 的 $\psi$ 类相交数都被确定。你能证明下面的闭公式吗：
% [Tutor Note: (2.61) 中必须把 $|\eps|<2$ 的那些项视为零，否则会出现 $(-1)!$、$(-2)!$ 这类无定义的阶乘。按此解释核对：$n=1,(d_1)=(1)$ 得 $\frac{1}{24}(1-0)=\frac{1}{24}$；$n=2,(1,1)$ 得 $\frac{1}{24}(2-1)=\frac{1}{24}$；$n=3,(1,1,1)$ 得 $\frac{1}{24}(6-4)=\frac{1}{12}$，而膨胀子方程给出 $\langle\tau_1^3\rangle_1=2\langle\tau_1^2\rangle_1=\frac{1}{12}$，一致。]
\begin{equation}
  \langle\tau_{d_1}\cdots\tau_{d_n}\rangle_1=\frac{1}{24}\left(\binom{n}{d_1,\dots,d_n}-\sum_{\eps_1,\dots,\eps_n\in\{0,1\}}\binom{n-|\eps|}{d_1-\eps_1,\dots,d_n-\eps_n}(|\eps|-2)!\right),
  \tag{2.61}
\end{equation}
其中 $|\eps|=\eps_1+\cdots+\eps_n$？
\end{exercise}

亏格 $0$ 的初始值 $\langle\tau_0^3\rangle_0=1$ 是平凡满足的，而亏格 $1$ 的情形 $\langle\tau_1\rangle_1=\frac{1}{24}$ 则相当不平凡。后者可以利用模空间 $\cM_{1,1}$ 的几何以及它与模形式的联系来计算。

\begin{exercise}[2.8]
利用以下事实证明 $\langle\tau_1\rangle_1=\frac{1}{24}$。
% [BOUNDARY: 原文 p.21 继续列举 Exercise 2.8 的 (2)(3)(4)…，属 E 块；本块在其第 (1) 项之后必须闭合 exercise/enumerate。E 块应从第 (2) 项续写（新开 enumerate 并 \setcounter{enumi}{1}），不要再开 exercise 环境，否则“习题 2.8”会印两遍。]
\begin{enumerate}[label=(\arabic*)]
  \item 对任意线丛 $\cL$，下述恒等式成立：$c_1(\cL)=\frac{1}{k}c_1(\cL^{\otimes k})$。
  % [Tutor Note: 该恒等式来自线丛 Chern 类的可加性 $c_1(\cL_1\otimes\cL_2)=c_1(\cL_1)+c_1(\cL_2)$（等价于 $c_1(\cL^{\otimes k})=k\,c_1(\cL)$）；需要 $k\ne 0$，取任意正整数即可。]
\end{enumerate}
\end{exercise}
